\documentclass[a4paper,10.5pt]{article}
\usepackage[utf8]{inputenc}
\usepackage[T1]{fontenc}
\usepackage[french]{babel}
\usepackage[left=1cm,right=1cm,top=.5cm,bottom=0cm]{geometry}
\usepackage{enumitem}
\usepackage{hyperref}
\usepackage{xcolor}
\usepackage{titlesec}
\usepackage{fontawesome5}
\usepackage{lmodern}
\usepackage[normalem]{ulem}

% --- Couleurs Compagnie des Signaux / MERMEC ---
\definecolor{primary}{RGB}{0, 80, 120}
\definecolor{secondary}{RGB}{80, 80, 80}

% --- Config Liens ---
\hypersetup{colorlinks=true, linkcolor=primary, filecolor=primary, urlcolor=primary}

% --- Titres ---
\titleformat{\section}{\large\bfseries\color{primary}\uppercase}{}{0em}{}[\titlerule]
\titlespacing{\section}{0pt}{8pt}{5pt}

% --- Commandes ---
\newcommand{\resumeItem}[1]{\item\small{#1}}

\begin{document}

% --- EN-TÊTE ---
\begin{center}
    {\Large \textbf{Loïc Aron MBASSI EWOLO}} \\ \vspace{4pt}
    \textbf{\large Alternance Tests Composants Logiciel Embarqué} \\
    \textit{\small Disponible dès septembre 2026 pour 12 mois} \\ \vspace{4pt}
    \small
    \faMapMarker\ Cergy, France (Mobile IDF) \quad
    \faPhone\ (+33) 7 59 52 33 98 \quad
    {\color{primary}\faEnvelope\ \href{mailto:loicaron.mbassiewolo@ensea.fr}{loicaron.mbassiewolo@ensea.fr}} \\ \vspace{2pt}
    {\color{primary}\faLinkedin\ \href{https://www.linkedin.com/in/loic-mbassi/}{linkedin.com/in/loic-mbassi}} \quad
    {\color{primary}\faGithub\ \href{https://github.com/Nameless0l}{github.com/Nameless0l}} \quad
\end{center}

\vspace{-6pt}

% --- PROFIL ---
\noindent\small{Développement d'un outil open source de génération automatique de code (Python, parsing avancé, +30 étoiles GitHub) et détection/isolation de défauts sur drone en environnement safety-critical (MATLAB/Simulink). Double diplôme ENSEA/Polytechnique Yaoundé en systèmes embarqués, signal et IA. Expérience en programmation système C/Linux (IPC, signaux, mémoire partagée) et en développement Python structuré (pipelines INRIA Saclay). Rigueur de programmation et documentation technique en français et anglais.}

% --- FORMATION ---
\section{Formation}
\begin{itemize}[leftmargin=*, label=\tiny$\bullet$, nosep]
    \item \textbf{ENSEA -- École Nationale Supérieure de l'Électronique et de ses Applications} \hfill \textit{2025 -- 2027} \\
    \small{Double diplôme d'Ingénieur -- Systèmes Embarqués, Traitement du Signal, FPGA/VHDL, IA.}
    \item \textbf{École Polytechnique de Yaoundé (1\textsuperscript{ère} école d'ingénieur CEMAC)} \hfill \textit{2021 -- 2025} \\
    \small{Diplôme d'Ingénieur de Conception (Master 2) : \textbf{Mathématiques Appliquées}, Génie Logiciel, Algorithmique.}
\end{itemize}

% --- COMPÉTENCES ---
\section{Compétences Techniques}
\begin{itemize}[leftmargin=*, nosep]
    \item \small{\textbf{Langages :} Python (pipelines, scripting, automatisation), C/C++ (embarqué, système), VBA, Java, PHP.}
    \item \small{\textbf{Embarqué \& Système :} STM32, Arduino, Raspberry Pi, Nvidia Jetson, Linux (IPC, SHM, signaux POSIX).}
    \item \small{\textbf{Tests \& Qualité :} Détection de défauts (FDI), contrôle tolérant aux pannes (FTC), MATLAB/Simulink, documentation technique.}
    \item \small{\textbf{Outils de développement :} Git, Docker, CI/CD, API REST, MySQL, Excel.}
    \item \small{\textbf{Génération de code \& Automatisation :} Parsing XML/UML, métaprogrammation, intégration LLM, scripting Python.}
    \item \small{\textbf{Langues :} Français (natif), Anglais (professionnel/technique -- rédaction scientifique INRIA).}
\end{itemize}

% --- EXPÉRIENCES / PROJETS ---
\section{Expériences \& Projets}
\begin{itemize}[leftmargin=*, label={-}]

\item \textbf{Fault Detection \& Fault-Tolerant Control -- Drone (Safety-Critical)} \hfill \textit{2025} \\
\textit{Projet d'option -- Contrôle automatique, ENSEA} \hfill \textit{MATLAB/Simulink}
\vspace{-2pt}
    \begin{itemize}[leftmargin=0.4cm, nosep, label=\tiny$\bullet$]
        \item \small{Détection et isolation de défauts (FDI) sur un système embarqué critique : perte d'actionneurs, biais et bruit capteurs.}
        \item \small{Contrôle tolérant aux pannes (FTC) avec reconfiguration de contrôleur et gain scheduling.}
        \item \small{Simulation de scénarios de pannes réalistes, évaluation comparative des performances.}
    \end{itemize}
\vspace{3pt}

\item \textbf{Laravel API Generator -- Outil Open Source (+30\faStarHalf\ GitHub)} \hfill \textit{2024 -- Présent} \\
\textit{Développement d'outil -- Automatisation \& Tests} \hfill \textit{PHP, Python, XML, Packagist}
\vspace{-2pt}
    \begin{itemize}[leftmargin=0.4cm, nosep, label=\tiny$\bullet$]
        \item \small{Conception d'un outil de génération automatique de code backend depuis des diagrammes UML/XML (parsing géométrique avancé).}
        \item \small{Correction de bugs, amélioration continue, vérification avec les utilisateurs, documentation technique complète.}
        \item \small{Déployé sur Packagist, maintenu en open source avec gestion des issues et retours utilisateurs.}
    \end{itemize}
\vspace{3pt}

\item \textbf{Stage Recherche -- INRIA Saclay, Équipe TRIBE} \hfill \textit{Avr 2026 -- Août 2026} \\
\textit{Plateau de Saclay / Institut Polytechnique de Paris} \hfill \textit{Python, Pandas, Scikit-learn, NLP}
\vspace{-2pt}
    \begin{itemize}[leftmargin=0.4cm, nosep, label=\tiny$\bullet$]
        \item \small{Pipeline Python d'analyse automatisée à grande échelle : code structuré, tests, documentation rigoureuse.}
        \item \small{Rédaction scientifique en anglais, collaboration multi-laboratoires.}
    \end{itemize}
\vspace{3pt}

\item \textbf{Challenge CoHoMa -- Robotique Autonome (ENSEA / Armée française)} \hfill \textit{2025 -- En cours} \\
\textit{Développement embarqué} \hfill \textit{C/C++, Nvidia Jetson, Docker, SLAM, Lidar 3D}
\vspace{-2pt}
    \begin{itemize}[leftmargin=0.4cm, nosep, label=\tiny$\bullet$]
        \item \small{Développement embarqué C/C++ sur Jetson : intégration capteurs (Lidar VLP16, caméra de profondeur), tests terrain.}
        \item \small{Containerisation Docker pour reproductibilité des environnements de test multi-plateforme.}
    \end{itemize}
\vspace{3pt}

\item \textbf{Simulation Flotte de Taxis -- Programmation Système C/Linux} \hfill \textit{2024} \\
\textit{Projet académique -- ENSPY} \hfill \textit{C, IPC (SHM, Signaux), OpenGL, Linux}
\vspace{-2pt}
    \begin{itemize}[leftmargin=0.4cm, nosep, label=\tiny$\bullet$]
        \item \small{Simulation multiprocessus en C : mémoire partagée, signaux POSIX, ordonnanceur FIFO, gestion manuelle mémoire.}
        \item \small{Rigueur de programmation bas niveau, synchronisation de processus concurrents.}
    \end{itemize}
\vspace{3pt}

\end{itemize}

% --- DISTINCTIONS ---
\section{Distinctions \& Leadership}
\begin{itemize}[leftmargin=*, label=\tiny$\bullet$, nosep]
    \item \small{\textbf{1\textsuperscript{er} Prix} CITS National 2024 (75 équipes) -- Optimisation ML/IoT collecte déchets.}
    \item \small{\textbf{Lead AI Cell} (ENSPY) : animation workshops techniques (Python, ML, architecture logicielle).}
    \item \small{\textbf{Certif.} Supervised ML (Stanford/DeepLearning.AI), Python for Data Science (IBM/Coursera).}
\end{itemize}

\end{document}
